\section{Related Work}
\label{sec:related_work}

\para{Strategic deception.} The strongest prior evidence for deliberate model misbehavior comes from trigger-conditioned settings. Sleeper Agents shows that models can preserve deceptive behavior through standard safety training and can even learn to hide that behavior better under adversarial training \citep{hubinger2024sleeper}. That work studies latent backdoors and persistent deception. Our setting is different. We do not claim strategic deception; instead, we measure an observable inference-time behavior in which the model produces compliant-looking output when abstention would be more appropriate.

\para{Truthfulness and calibrated factual answering.} \truthfulqa measures whether models reproduce common human falsehoods rather than truthful answers \citep{lin2021truthfulqa}. \simpleqa asks a complementary question: does a model ``know what it knows,'' and can it leave a question unanswered when needed \citep{wei2024simpleqa}. These benchmarks are directly relevant because they separate correct answers from plausible but wrong answers and, in the case of \simpleqa, explicitly allow non-attempts. We build on them by keeping the question set fixed while manipulating answer pressure, which makes the abstention-versus-bluffing tradeoff directly measurable.

\para{Hallucination benchmarks and checkers.} Hallucination has been studied through benchmark construction and verifier design. HaluEval provides a large benchmark for hallucination recognition across multiple tasks \citep{li2023halueval}. RefChecker introduces claim-triplet verification for fine-grained hallucination detection \citep{hu2024refchecker}. \halogen goes further by releasing prompts, verifiers, and an error taxonomy across multiple domains \citep{ravichander2025halogen}. We use \halogen's false-presupposition prompts because they offer a clean impossible-request setting where abstention is not only acceptable but correct.

\para{Abstention under uncertainty.} AbstentionBench argues that abstention remains an unsolved reliability problem, especially for reasoning-tuned models on unanswerable questions \citep{kirichenko2025abstentionbench}. Our work complements that line by experimentally varying prompt incentives. Instead of only measuring whether models abstain when they should, we test whether removing the abstention option causes a predictable shift into pseudo-compliance or fabrication.

Taken together, these lines of work suggest that deception, hallucination, truthfulness failure, and abstention failure are related but usually evaluated in isolation. Our contribution is to study them through a shared operational lens: whether answer pressure makes models choose compliant output over epistemically honest behavior.
